\section{Introduction}

In late 2023, FunSearch demonstrated that a large language model
inside an evolutionary loop can produce genuinely new mathematics
\cite{romera2024funsearch}; AlphaEvolve scaled the recipe to dozens
of problems and whole codebases \cite{novikov2025alphaevolve}. These
systems are typically reported as monoliths: a model, a mutation
scheme, an evaluator, a database --- and a result. When the result is
good, it is hard to say which ingredient earned it. Practitioners
inherit an expensive ambiguity: should the next budget go to more
iterations, a bigger population, a stronger model --- or to a
different \emph{mode} of the same model?

This paper isolates one ingredient experimentally: the model's
extended reasoning mode (``thinking''), in which the model produces a
long private chain of thought before its answer. Our vehicle is a
campaign on covering designs --- minimum collections of $k$-element
blocks covering all $t$-subsets of a $v$-element universe --- chosen
for three properties. First, \emph{verifier asymmetry}: constructing
good coverings has occupied specialists for decades, but verifying a
candidate is exact integer counting, which suits our harness's
governing rule that nothing an evolved program says about itself is
ever trusted. Second, the domain has a canonical, curated frontier
(the La Jolla Covering Repository and its live successor
\cite{ljcr}) whose values on our target cells have stood for 15--30
years --- a stagnant frontier against which even exact ties are
informative. Third, to our knowledge no system in the FunSearch
lineage has touched it.

The campaign itself furnished the experiment. Cheap ensemble
evolution improved $C(32,8,4)$ from a 1{,}258-block seed baseline to
977 blocks within five mutations, then plateaued. We spent 75
iterations establishing that the plateau was real rather than an
instrument artifact: we repaired a mutation-waste pathology (48\% of
calls were being rejected before evaluation), verified the fix (waste
fell to 8\%), and escalated to the conventional remedy --- full
rewrites by the flagship model --- which produced children clustering
to the \emph{exact} plateau value from four independent rewrites.
Then we enabled extended reasoning on the same model, same prompt,
same frozen population. The first improving iteration rewrote the
solver around a general finite-geometry construction --- all
$d$-flats of $\mathrm{AG}(m,p)$ for any prime-power universe --- whose
block count at $(32,8,4)$ is exactly 620: the repository's best-known
value, dated 1996, and 357 blocks (36\%) below the plateau. A
controlled repetition experiment from the frozen state confirmed the
effect is categorical, not luck: three of three reasoning-enabled
slices rediscovered the construction independently; zero of three
reasoning-disabled slices left the plateau band ($p = 0.05$,
one-sided Fisher; 55 disabled full-rewrite attempts in total never
produced a structural change).

We then asked what the discovery is worth beyond its cell. Under a
benchmark protocol frozen before execution --- 29 cells spanning
proven-optimum gates, affine-family holdouts, and cells where the
affine construction cannot apply --- the single evolved program
matched the best-known value in 22 cells with zero per-cell
adaptation, including $C(49,8,2)$, where the pure construction is
seven blocks short and the evolved local search closes the gap, and
four non-prime-power cells served purely by the evolved search
machinery. All ties were re-verified by a standalone verifier
sharing no code with the harness. Because ``the model memorized the
tables'' is the natural objection, we also ran dissociation probes:
the deployed model, asked directly, declines to state the value of
$C(32,8,4)$; and asked to emit a covering as plain text with a
64k-token reasoning budget but no code, it produces zero blocks. The
capability that matters is not recall but the
reasoning-to-code-to-execution pathway.

Our contributions:

\begin{description}
\item[C1 --- A controlled reasoning-mode experiment in an evolutionary
discovery loop] (Section~\ref{sec:experiment}): to our knowledge the
first, showing a categorical 3/3-vs-0/3 breakthrough separation with
a 36\% objective effect, against a backdrop of mixed
reasoning-ablation results in other task families.
\item[C2 --- Verified rediscovery at a stagnant frontier]
(Sections~\ref{sec:campaign}, \ref{sec:benchmark}): 22 of 29
pre-registered cells tied to best-known values by one evolved
program, every tie independently verified; covering designs enter the
LLM-evolutionary literature.
\item[C3 --- An exactness-first, problem-agnostic harness]
(Sections~\ref{sec:harness}--\ref{sec:instrument}): plugin contract
with fitness/cost separation, multi-seed fitness, work-counter
determinism, solution archiving, and leak-audited prompts --- each
mechanism motivated by a documented failure; plus engineering
findings (mutation-waste dynamics, an ensemble-seeding pathology)
reported for reuse.
\item[C4 --- A contamination-dissociation methodology]
(Section~\ref{sec:contamination}): value-recall, direct-artifact, and
code-pathway probes that calibrate what ``rediscovery'' means when
the literature is in the weights.
\end{description}

Section~\ref{sec:related} reviews related work;
Section~\ref{sec:limitations} details limitations --- foremost that
all positive results are ties, not records, and that the causal claim
is scoped to one model, one plateau, one domain. Everything needed to
re-verify or extend this work --- harness, certificates, protocols,
raw logs --- is released as open source.
